% Source PDF: source/普通高考/2012/2012浙江文.pdf, page 1, question 13.
% Elements: rounded start/end terminals, three process rectangles, one decision
% diamond, input/output parallelogram, and the loop-back connector. All borders
% and connectors are solid; no dashed segments or hidden edges are present.
% Node -> directed-edge list: start -> (T=1,i=1) -> (T=T/i) -> (i=i+1)
% -> (i>5); yes -> output T -> finish; no -> loop entry immediately before
% T=T/i. The loop entry carries one left arrowhead and the main stem carries
% one downward arrowhead into T=T/i (no overlapping duplicate arrows).
\begin{tikzpicture}[
  >=stealth,
  x=1cm,y=1cm,
  line width=0.45pt,
  line cap=round,line join=round,
  font=\normalsize,
  flowbox/.style={draw,anchor=center,align=center,inner xsep=4pt,inner ysep=2pt,
    execute at begin node={\setlength{\parindent}{0pt}}},
  every node/.style={anchor=center,align=center,
    execute at begin node={\setlength{\parindent}{0pt}}},
  terminal/.style={flowbox,rounded corners=.17cm,minimum width=1.35cm,minimum height=.72cm},
  process/.style={flowbox,minimum height=.76cm},
  io/.style={flowbox,shape=trapezium,trapezium left angle=72,trapezium right angle=108,
    minimum width=2.20cm,minimum height=.78cm},
  decision/.style={flowbox,diamond,aspect=2.7,minimum width=3.85cm,minimum height=1.00cm},
]
  % Safety margin keeps the rounded end terminal and loop line inside the bbox.
  \path[use as bounding box] (-2.05,-8.35) rectangle (2.95,0.40);

  \node[terminal] (start) at (0,0) {开始};
  \node[process,minimum width=2.85cm] (init) at (0,-1.18) {$T=1,i=1$};
  \node[process,minimum width=1.95cm,minimum height=1.00cm] (divide) at (0,-2.68) {$T=\dfrac{T}{i}$};
  \node[process,minimum width=2.15cm] (increment) at (0,-3.90) {$i=i+1$};
  \node[decision] (test) at (0,-5.18) {$i>5?$};
  \node[io,minimum width=2.25cm] (output) at (0,-6.50) {输出 $T$};
  \node[terminal] (finish) at (0,-7.78) {结束};
  % Keep the return connector midway between init and divide with >=2.5 mm clearance.
  \coordinate (merge-level) at (0,-1.82);

  \draw[->] (start.south) -- (init.north);
  \draw (init.south) -- (merge-level);
  \draw[->] (merge-level) -- (divide.north);
  \draw[->] (divide.south) -- (increment.north);
  \draw[->] (increment.south) -- (test.north);
  \draw[->] (test.south) -- node[right=2pt] {是} (output.north);
  \draw[->] (output.south) -- (finish.north);

  % False branch exits right, rises, and points left into the T update entry.
  \coordinate (loop-right) at (2.45,-5.18);
  \draw (test.east) -- node[above=2pt] {否} (loop-right) -- (2.45,-1.82);
  % The return arrow ends at the degree-3 merge; the main stem continues with
  % its independent downward arrow into the division box.
  \draw[->] (2.45,-1.82) -- (0.75,-1.82) -- (merge-level);
\end{tikzpicture}
